\newsection{4} {Linear Higher Order Equations}
{Variation of Parameters for Higher Order  Equations}
\currentpdfbookmark{Section 9.4 Variation of Parameters for Higher Order  Equations}{section:9.4}
\vskip14pt

\renewcommand{\thissection}{\sectiontitle
{\,VARIATION OF PARAMETERS  FOR HIGHER ORDER  EQUATIONS}}
\thissection

\noindent
\boxit{Derivation of the method}

\noindent
We assume throughout this section that the nonhomogeneous linear
equation
\begin{equation} \label{eq:9.4.1}
P_0(x)y^{(n)}+P_1(x)y^{(n-1)}+\cdots+P_n(x)y=F(x)
\end{equation}
is normal on an interval $(a,b)$. We'll abbreviate this equation as
$Ly=F$, where
$$
Ly=P_0(x)y^{(n)}+P_1(x)y^{(n-1)}+\cdots+P_n(x)y.
$$
When we speak of solutions of this equation and its complementary
equation $Ly=0$, we mean solutions on $(a,b)$. We'll show how to use
the method of variation of parameters to find a particular solution of
$Ly=F$, provided that we know a fundamental set of solutions
$\{y_1,y_2,\dots,y_n\}$ of $Ly=0$.

We seek a particular solution of $Ly=F$ in the form
\begin{equation} \label{eq:9.4.2}
y_p=u_1y_1+u_2y_2+\cdots+u_ny_n
\end{equation}
where $\{y_1,y_2,\dots,y_n\}$ is a known fundamental set of solutions
of the complementary equation
$$
P_0(x)y^{(n)}+P_1(x)y^{(n-1)}+\cdots+P_n(x)y=0
$$
and $u_1$, $u_2$, \dots, $u_n$ are functions to be determined. We begin by
imposing the following $n-1$ conditions on $u_1,u_2,\dots,u_n$:
\begin{equation} \label{eq:9.4.3}
\begin{array}{rcl}
u'_1y_1+u'_2y_2+&\cdots&+u'_ny_n=0 \\
u'_1y'_1+u'_2y'_2+&\cdots&+u'_ny'_n=0 \\
\phantom{u'_1y^{(n_1)}+u'_2y_2^{(n-1)}}&\vdots&
\phantom{\cdots+u'_ny^{(n-1)}_n=q} \\
u'_1y_1^{(n-2)}+u'_2y^{(n-2)}_2+&\cdots&+u'_ny^{(n-2)}_n
=0.
\end{array}
\end{equation}
These conditions lead to simple formulas for the first $n-1$
derivatives of $y_p$:
\begin{equation} \label{eq:9.4.4}
y^{(r)}_p=u_1y^{(r)}_1+u_2y_2^{(r)}\cdots+u_ny^{(r)}_n,\ 0
\le r \le n-1.
\end{equation}
These formulas are easy to remember, since they look as though we
obtained them by differentiating \eqref{eq:9.4.2} $n-1$ times while
treating $u_1$, $u_2$, \dots, $u_n$ as constants. To see that
\eqref{eq:9.4.3}
implies \eqref{eq:9.4.4}, we first differentiate \eqref{eq:9.4.2} to obtain
$$
y_p'=u_1y_1'+u_2y_2'+\cdots+u_ny_n'+u_1'y_1+u_2'y_2+\cdots+u_n'y_n,
$$
which  reduces to
$$
y_p'=u_1y_1'+u_2y_2'+\cdots+u_ny_n'
$$
because of the first equation in \eqref{eq:9.4.3}. Differentiating this
yields
$$
y_p''=u_1y_1''+u_2y_2''+\cdots+u_ny_n''+u_1'y_1'+u_2'y_2'+\cdots+u_n'y_n',
$$
which reduces to
$$
y_p''=u_1y_1''+u_2y_2''+\cdots+u_ny_n''
$$
because of  the second equation in \eqref{eq:9.4.3}.
Continuing in this way yields \eqref{eq:9.4.4}.

The last equation in \eqref{eq:9.4.4} is
$$
y_p^{(n-1)}=u_1y_1^{(n-1)}+u_2y_2^{(n-1)}+\cdots+u_ny_n^{(n-1)}.
$$
Differentiating this yields
$$
y_p^{(n)}=u_1y_1^{(n)}+u_2y_2^{(n)}+\cdots+u_ny_n^{(n)}+
u_1'y_1^{(n-1)}+u_2'y_2^{(n-1)}+\cdots+u_n'y_n^{(n-1)}.
$$
Substituting this and  \eqref{eq:9.4.4} into \eqref{eq:9.4.1}
yields
$$
u_1Ly_1+u_2Ly_2+\cdots+u_nLy_n+P_0(x)\left(
u_1'y_1^{(n-1)}+u_2'y_2^{(n-1)}+\cdots+u_n'y_n^{(n-1)}\right)=F(x).
$$
Since  $Ly_i=0$ $(1 \le i \le n)$, this reduces to
$$
u_1'y_1^{(n-1)}+u_2'y_2^{(n-1)}+\cdots+u_n'y_n^{(n-1)}={F(x)\over P_0(x)}.
$$
Combining this equation with \eqref{eq:9.4.3} shows that
$$
y_p=u_1y_1+u_2y_2+\cdots+u_ny_n
$$
is a solution of \eqref{eq:9.4.1} if
$$
\begin{array}{rcl}
u'_1y_1+u'_2y_2+&\cdots&+u'_ny_n=0 \\
u'_1y'_1+u'_2y'_2+&\cdots&+u'_ny'_n=0 \\
\phantom{u'_1y^{(n_1)}+u'_2y_2^{(n-1)}}&\vdots&
\phantom{\cdots+u'_ny^{(n-1)}_n=q} \\
u'_1y_1^{(n-2)}+u'_2y^{(n-2)}_2+&\cdots&+u'_ny^{(n-2)}_n
=0 \\
u'_1y^{(n-1)}_1+u'_2y^{(n-1)}_2+&\cdots&+u'_n
y^{(n-1)}_n=F/P_0,
\end{array}
$$
which can be written in matrix form as
\begin{equation} \label{eq:9.4.5}
\left[\begin{array}{cccc}
y_1&y_2&\cdots&y_n \\[2\jot]
y'_1&y'_2&\cdots&y_n'\\[2\jot]
\vdots&\vdots&\ddots&\vdots\\[2\jot]
y_1^{(n-2)}&y_2^{(n-2)}&\cdots&y_n^{(n-2)}\\[2\jot]
y_1^{(n-1)}&y_2^{(n-1)}&\cdots&y_n^{(n-1)}
\end{array} \right]
\left[\begin{array}{c}u_1'\\u_2'\\\vdots\\u_{n-1}'\\u_n'\end{array}
\right]=
\left[\begin{array}{c}0\\0\\ \vdots\\0\\F/ P_0\end{array}\right].
\end{equation}
The determinant of this system  is the Wronskian $W$ of
the fundamental set of solutions  $\{y_1,y_2,\dots,y_n\}$,
which has no zeros on $(a,b)$, by Theorem~\ref{thmtype:9.1.4}.
Solving
\eqref{eq:9.4.5} by Cramer's rule  yields
\begin{equation} \label{eq:9.4.6}
u'_j=(-1)^{n-j}{FW_j\over P_0W},\quad 1\le j\le n,
\end{equation}
 where $W_j$ is the Wronskian of the set of
functions obtained by deleting $y_j$ from  $\{y_1,y_2,\dots,y_n\}$
 and keeping the remaining functions in the
same order. Equivalently, $W_j$ is the determinant obtained by
deleting the last row and $j$-th column of $W$.

Having obtained $u_1'$, $u_2'$$, \dots, $$u_n'$, we can integrate  to
obtain $u_1,\,u_2,\dots,u_n$. As in Section~5.7, we
take the constants
of integration to be zero, and we drop any linear combination of
$\{y_1,y_2,\dots,y_n\}$ that may appear in $y_p$.

\medskip

\color{blue}
\remark{ For efficiency, it's best to compute $W_1$, $W_2$, \dots, $W_n$
first, and then compute $W$ by expanding in cofactors of the last row;
thus,
$$
W=\sum_{j=1}^n(-1)^{n-j}y_j^{(n-1)}W_j.
$$}
\color{black}

\boxit{Third Order Equations}

\noindent
If $n=3$, then
$$
W=\left| \begin{array}{ccc}
y_1&y_2&y_3 \\[2\jot]
y'_1&y'_2&y'_3 \\[2\jot]
y''_1&y''_2&y''_3 \end{array} \right|.
$$
Therefore
$$
W_1=\left| \begin{array}{cc}
y_2&y_3 \\[2\jot]
y'_2&y'_3 \end{array} \right|,
\quad W_2=\left| \begin{array}{cc}
y_1&y_3 \\[2\jot]
 y'_1&y'_3 \end{array} \right|,
\quad W_3=\left| \begin{array}{cc}
y_1&y_2 \\[2\jot]
y'_1&y'_2 \end{array} \right|,
$$
and  \eqref{eq:9.4.6}  becomes
\begin{equation} \label{eq:9.4.7}
u'_1={FW_1\over P_0W},\quad u'_2=-{FW_2\over P_0W},\quad
u'_3={FW_3\over P_0W}.
\end{equation}

\begin{example}\label{example:9.4.1} \rm
 Find a particular solution of
\begin{equation} \label{eq:9.4.8}
xy'''-y''-xy'+y=8x^2e^x,
\end{equation}
given that  $y_1=x$, $y_2=e^x$, and $y_3=e^{-x}$ form a fundamental
set of solutions of the complementary equation. Then find the general
solution of \eqref{eq:9.4.8}.
 \end{example}

\solution
We seek a particular solution of \eqref{eq:9.4.8} of the form
$$
y_p=u_1x+u_2e^x+u_3e^{-x}.
$$
The Wronskian of  $\{y_1,y_2,y_3\}$ is
$$
W(x)=\left| \begin{array}{ccr}
 x&e^x&e^{-x} \\ 1&e^x&-e^{-x} \\ 0&e^x&e^{-x}
\end{array} \right|,
$$
so
\begin{eqnarray*}
W_1&=&
\left| \begin{array}{cr}
e^x&e^{-x}\\ e^x&-e^{-x}
\end{array} \right|=-2,\\[10pt]
W_2&=&
\left| \begin{array}{cr}
x&e^{-x}\\1&-e^{-x}
\end{array} \right|=-e^{-x}(x+1),\\
W_3&=&
\left| \begin{array}{cc}
x&e^x\\1&e^x
\end{array} \right|=e^x(x-1).
\end{eqnarray*}
Expanding $W$ by cofactors of the last row yields
$$
W=0W_1-e^x W_2+e^{-x}W_3=0(-2)-e^x\left(-e^{-x}(x+1)\right)
+e^{-x}e^x(x-1)=2x.
$$
 Since $F(x)=8x^2e^x$ and $P_0(x)=x$,
$$
{F\over P_0W}={8x^2e^x\over x\cdot 2x}=4e^x.
$$
Therefore, from \eqref{eq:9.4.7}
\begin{eqnarray*}
u'_1&=&\phantom{-} 4e^xW_1=\phantom{-}4e^x(-2)=-8e^x,\\
u_2'&=&-4e^xW_2=-4e^x\left(-e^{-x}(x+1)\right)=4(x+1),\\
u_3'&=&\phantom{-}4e^xW_3=\phantom{-}4e^x\left(e^x(x-1)\right)=4e^{2x}(x-1).
\end{eqnarray*}
 Integrating and taking the constants of
integration to be zero yields
$$
u_1=-8e^x,\quad u_2=2(x+1)^2, u_3=e^{2x}(2x-3).
$$
Hence,
\begin{eqnarray*}
y_p&=&u_1y_1+u_2y_2+u_3y_3\\
&=&(-8e^x)x+e^x(2(x+1)^2)+e^{-x}\left(e^{2x}(2x-3)\right)
\\&=&e^x(2x^2-2x-1).
\end{eqnarray*}
Since $-e^x$ is a solution of the complementary equation, we redefine
$$
y_p=2xe^x(x-1).
$$
 Therefore the general solution of \eqref{eq:9.4.8} is
$$
y=2xe^x(x-1)+c_1x+c_2e^x+c_3e^{-x}.
$$

\boxit{Fourth Order Equations}

\noindent
If $n=4$, then
$$
W=\left| \begin{array}{cccc}
y_1&y_2&y_3&y_4 \\[2\jot]
y'_1&y'_2&y'_3&y_4' \\[2\jot]
y''_1&y''_2&y''_3&y_4''\\[2\jot]
y'''_1&y'''_2&y'''_3&y_4'''
 \end{array} \right|,
$$
Therefore
$$
W_1=\left| \begin{array}{ccc}
y_2&y_3&y_4 \\[2\jot]
y'_2&y'_3&y_4'\\[2\jot]
y''_2&y''_3&y_4''
\end{array} \right|,
\quad W_2=\left| \begin{array}{ccc}
y_1&y_3&y_4 \\[2\jot]
y'_1&y'_3&y_4'\\[2\jot]
y''_1&y''_3&y_4''
 \end{array} \right|,
$$
\vskip5pt
$$
 W_3=\left| \begin{array}{ccc}
y_1&y_2&y_4 \\[2\jot]
y'_1&y'_2&y_4'\\[2\jot]
y''_1&y''_2&y_4''
 \end{array} \right|,\quad
 W_4=\left| \begin{array}{ccc}
y_1&y_2&y_3 \\[2\jot]
y_1'&y'_2&y_3'\\[2\jot]
y_1''&y''_2&y_3''
 \end{array} \right|,
$$
and  \eqref{eq:9.4.6}  becomes
\begin{equation} \label{eq:9.4.9}
u'_1=-{FW_1\over P_0W},\quad u'_2={FW_2\over P_0W},\quad
u'_3=-{FW_3\over P_0W},\quad u'_4={FW_4\over P_0W}.
\end{equation}

\begin{example}\label{example:9.4.2} \rm
 Find a particular solution of
\begin{equation} \label{eq:9.4.10}
x^4y^{(4)}+6x^3y'''+2x^2y''-4xy'+4y=12x^2,
\end{equation}
given that  $y_1=x$, $y_2=x^2$, $y_3=1/x$ and $y_4=1/x^2$ form a
fundamental set of solutions of the complementary equation. Then find
the general solution of \eqref{eq:9.4.10} on $(-\infty,0)$ and
$(0,\infty)$.
 \end{example}

\solution
We seek a particular solution of \eqref{eq:9.4.10} of the form
$$
y_p=u_1x+u_2x^2+{u_3\over x}+{u_4\over x^2}.
$$
The Wronskian of  $\{y_1,y_2,y_3,y_4\}$ is
$$
W(x)=\left| \begin{array}{cccr}
 x&x^2&1/x&-1/x^2 \\[5pt] 1&2x&-1/x^2&-2/x^3 \\[5pt] 0
&2&2/x^3&6/x^4\\[5pt]
0&0&-6/x^4&-24/x^5
\end{array} \right|,
$$
so
\begin{eqnarray*}
W_1&=&
\left| \begin{array}{ccr}
x^2&1/x&1/x^2\\[5pt]2x&-1/x^2&-2/x^3\\[5pt]
2&2/x^3&6/x^4
\end{array} \right|=-{12\over x^4},\\[10pt]
W_2&=&
\left| \begin{array}{ccr}
x&1/x&1/x^2\\[5pt]1&-1/x^2&-2/x^3\\[5pt]
0&2/x^3&6/x^4
\end{array} \right|=-{6\over x^5},\\[10pt]
W_3&=&
\left| \begin{array}{ccc}
x&x^2&1/x^2\\[5pt]1&2x&-2/x^3\\[5pt]
0&2&6/x^4
\end{array} \right|={12\over x^2},  \\
W_4&=&
\left| \begin{array}{ccc}
x&x^2&1/x\\[5pt]1&2x&-1/x^2\\[5pt]
0&2&2/x^3
\end{array} \right|={6\over x}.
\end{eqnarray*}
Expanding $W$ by cofactors of the last row yields
\begin{eqnarray*}
W&=&-0W_1+0 W_2-\left(-{6\over
x^4}\right)W_3+\left(-{24\over
x^5}\right)W_4\\
&=&{6\over x^4}{12\over x^2}-{24\over x^5}{6\over x}=-{72\over x^6}.
\end{eqnarray*}
 Since $F(x)=12x^2$ and $P_0(x)=x^4$,
$$
{F\over P_0W}={12x^2\over x^4}\left(-{x^6\over72}\right)=-{x^4\over 6}.
$$
Therefore, from \eqref{eq:9.4.9},
\begin{eqnarray*}
u'_1&=&-\left(-{x^4\over6}\right)W_1={x^4\over6}\left(-{12\over
x^4}\right)=-2,\\[5pt]
 u_2'&=&\phantom{-}-{x^4\over6}W_2=-{x^4\over6}\left(-{6\over x^5}\right)
={1\over x},\\[5pt]
u_3'&=&-\left(-{x^4\over6}\right)W_3={x^4\over6}{12\over x^2}=2x^2,\\
  u_4'&=&\phantom{-}-{x^4\over6}W_4=-{x^4\over6}{6\over x}=-x^3.
\end{eqnarray*}
 Integrating these and taking the constants of
integration to be zero yields
$$
u_1=-2x,\quad u_2=\ln|x|,\quad u_3={2x^3\over3},
u_4=-{x^4\over4}.
$$
Hence,
\begin{eqnarray*}
y_p&=&u_1y_1+u_2y_2+u_3y_3+u_4y_4\\
&=&(-2x)x+(\ln|x|)x^2+{2x^3\over3}{1\over x}+\left(-{x^4\over4}\right)
{1\over x^2} \\&=&x^2\ln|x|-{19x^2\over12}.
\end{eqnarray*}
Since $-19x^2/12$ is a solution of the complementary equation, we redefine
$$
y_p=x^2\ln|x|.
$$
 Therefore
$$
y=x^2\ln|x|+c_1x+c_2x^2+{c_3\over x}+{c_4\over x^2}
$$
is the general solution of \eqref{eq:9.4.10} on $(-\infty,0)$ and $(0,\infty)$.

\enlargethispage{1in}
\exercises
In Exercises~\ref{exer:9.4.1}--\ref{exer:9.4.21} find a particular
solution,
given the fundamental set of solutions of the complementary equation.

\begin{exerciselist}
\item\label{exer:9.4.1}
$x^3y'''-x^2(x+3)y''+2x(x+3)y'-2(x+3)y=-4x^4$;\quad
$\{x,\,x^2,\,xe^x\}$
\newpage

\item\label{exer:9.4.2}
$y'''+6xy''+(6+12x^2)y'+(12x+8x^3)y=x^{1/2}e^{-x^2}$;\quad
 $\{e^{-x^2},\,xe^{-x^2},\,x^2e^{-x^2}\}$

\item\label{exer:9.4.3} $x^3y'''-3x^2y''+6xy'-6y=2x$;
 \quad  $\{x,x^2,x^3\}$

\item\label{exer:9.4.4}
$x^2y'''+2xy''-(x^2+2)y'=2x^2$;\quad   $\{1,\,e^x/x,\,e^{-x}/x\}$

\item\label{exer:9.4.5}
$x^3y'''-3x^2(x+1)y''+3x(x^2+2x+2)y'-(x^3+3x^2+6x+6)y=x^4e^{-3x}$;\quad

$\{xe^x,\,x^2e^x,\,x^3e^x\}$

\item\label{exer:9.4.6}
$x(x^2-2)y'''+(x^2-6)y''+x(2-x^2)y'+(6-x^2)y=2(x^2-2)^2$; \quad
 $\{e^x,\,e^{-x},\,1/x\}$

\item\label{exer:9.4.7}
$xy'''-(x-3)y''-(x+2)y'+(x-1)y=-4e^{-x}$; \quad
$\{e^x,\,e^x/x,\,e^{-x}/x\}$

\item\label{exer:9.4.8}
$4x^3y'''+4x^2y''-5xy'+2y=30x^2$; \quad  $\{\sqrt x,\,1/\sqrt x,\,x^2\}$

\item\label{exer:9.4.9}
$x(x^2-1)y'''+(5x^2+1)y''+2xy'-2y=12x^2$; \quad
$\{x,\,1/(x-1),\,1/(x+1)\}$

\item\label{exer:9.4.10}
$x(1-x)y'''+(x^2-3x+3)y''+xy'-y=2(x-1)^2$; \quad
$\{x,\,1/x,e^x/x\}$

\item\label{exer:9.4.11}
$x^3y'''+x^2y''-2xy'+2y=x^2$;  \quad  $\{x,\,x^2,\,1/x\}$

\item\label{exer:9.4.12} $xy'''-y''-xy'+y=x^2$; \quad
$\{x,\,e^x,\,e^{-x}\}$

\item\label{exer:9.4.13}
$xy^{(4)}+4y'''=6 \ln |x|$; \quad   $\{1,\,x,\,x^2,\,1/x\}$

\item\label{exer:9.4.14}
$16x^4y^{(4)}+96x^3y'''+72x^2y''-24xy'+9y=96x^{5/2}$; \quad
$\{\sqrt x,\,1/\sqrt x,\,x^{3/2},\,x^{-3/2}\}$

\item\label{exer:9.4.15}
$x(x^2-6)y^{(4)}+2(x^2-12)y'''+x(6-x^2)y''+2(12-x^2)y'=2(x^2-6)^2$; \quad

 $\{1,\,1/x,\,e^x,\,e^{-x}\}$

\item\label{exer:9.4.16}
$x^4y^{(4)}-4x^3y'''+12x^2y''-24xy'+24y=x^4$; \quad
$\{x,\,x^2,\,x^3,\,x^4\}$

\item\label{exer:9.4.17}
$x^4y^{(4)}-4x^3y'''+2x^2(6-x^2)y''+4x(x^2-6)y'+(x^4-4x^2+24)y=4x^5e^x$;
\quad

 $\{xe^x,\,x^2e^x,\,xe^{-x},\,x^2e^{-x}\}$

\item\label{exer:9.4.18}
$x^4y^{(4)}+6x^3y'''+2x^2y''-4xy'+4y=12x^2$;   \quad
$\{x,x^2,1/x,1/x^2\}$

\item\label{exer:9.4.19}
$xy^{(4)}+4y'''-2xy''-4y'+xy=4e^x$; \quad
$\{e^x,\,e^{-x},\,e^x/x,\,e^{-x}/x\}$

\item\label{exer:9.4.20}
$xy^{(4)}+(4-6x)y'''+(13x-18)y''+(26-12x)y'+(4x-12)y=3e^x$; \quad
 $\{e^x,\,e^{2x},\,e^x/x,\,e^{2x}/x\}$

\item\label{exer:9.4.21}
$x^4y^{(4)}-4x^3y'''+x^2(12-x^2)y''+2x(x^2-12)y'+2(12-x^2)y=2x^5$;
 $\{x,\,x^2,\,xe^x,\,xe^{-x}\}$

\exercisetext{In Exercises~\ref{exer:9.4.22}--\ref{exer:9.4.33} solve the
initial value problem, given the fundamental set of solutions of the
complementary equation. Where indicated by \CGex, graph the solution.}

\item\label{exer:9.4.22} \CGex
$x^3y'''-2x^2y''+3xy'-3y=4x, \quad  y(1)=4,\quad y'(1)=4,
\quad y''(1)=2$; \quad     $\{x,\,x^3,\,x \ln x\}$

\item\label{exer:9.4.23}
$x^3y'''-5x^2y''+14xy'-18y=x^3, \quad  y(1)=0,\quad y'(1)=1,\quad
y''(1)=7$;
\quad $\{x^2,\, x^3,\,x^3 \ln x\}$

\item\label{exer:9.4.24}
$\dst{(5-6x)y'''+(12x-4)y''+(6x-23)y'+(22-12x)y=-(6x-5)^2e^x}$

$\dst{y(0)=-4, \quad  y'(0)=-{3\over2},\quad y''(0)=-19}$;\quad
 $\{e^x,\,e^{2x},\,xe^{-x} \}$

\item\label{exer:9.4.25}
$x^3y'''-6x^2y''+16xy'-16y=9x^4, \quad  y(1)=2,\quad y'(1)=1,\quad y''(1)=5$;

 $\{x,\,x^4,\,x^4 \ln |x|\}$

\item\label{exer:9.4.26} \CGex
$(x^2-2x+2)y'''-x^2y''+2xy'-2y=(x^2-2x+2)^2, \quad  y(0)=0,\quad y'(0)=5$,

$y''(0)=0$;\quad       $\{x,\,x^2,\,e^x\}$

\item\label{exer:9.4.27}
$x^3y'''+x^2y''-2xy'+2y=x(x+1), \quad  y(-1)=-6,\quad y'(-1)=\dst{43\over6},\quad y''(-1)=
-\dst{5\over2}$;\quad

  $\{x,\,x^2,\,1/x\}$
\newpage

\item\label{exer:9.4.28}
$(3x-1)y'''-(12x-1)y''+9(x+1)y'-9y=2e^x(3x-1)^2, \quad  y(0)=\dst{3\over4}$,

$y'(0)=\dst{5\over4}, \quad  y''(0)=\dst{1\over4}$;\quad
$\{x+1,\,e^x,\,e^{3x}\}$

\item\label{exer:9.4.29}  \CGex
$(x^2-2)y'''-2xy''+(2-x^2)y'+2xy=2(x^2-2)^2, \quad  y(0)=1,\quad y'(0)=-5$,

 $y''(0)=5$;\quad      $\{x^2,\,e^x,\,e^{-x}\}$

\item\label{exer:9.4.30} \CGex
$x^4y^{(4)}+3x^3y'''-x^2y''+2xy'-2y=9x^2, \quad  y(1)=-7,\quad y'(1)=
-11,\quad y''(1)=-5$,\;
$y'''(1)=6; \quad \{x,\,x^2,\,1/x,\,x\ln x\}$

\item\label{exer:9.4.31}
$(2x-1)y^{(4)}-4xy'''+(5-2x)y''+4xy'-4y=6(2x-1)^2, \quad  y(0)=\dst{55\over4},
\quad y'(0)=0$,

$y''(0)=13, \quad  y'''(0)=1$;\quad    $\{x,\,e^x,\,e^{-x},\,e^{2x}\}$

\item\label{exer:9.4.32}
$4x^4y^{(4)}+24x^3y'''+23x^2y''-xy'+y=6x$,\quad $y(1)=2,\quad
y'(1)=0,\quad y''(1)=4,\quad y'''(1)=-\dst{37\over4}$;
\quad $\{x,\sqrt x,1/x,1/\sqrt x\}$

\item\label{exer:9.4.33}
$x^4y^{{(4)}}+5x^3y'''-3x^2y''-6xy'+6y=40x^3,
\quad  y(-1)=-1, \;  y'(-1)=-7$,

$y''(-1)=-1,\quad y'''(-1)=-31$;\quad
$\{x,\, x^3,\,1/x,\,1/x^2\}$

\item\label{exer:9.4.34}
Suppose the equation
$$
P_0(x)y^{(n)}+P_1(x)y^{(n-1)}+\cdots+P_n(x)y=F(x)
\eqno{\rm (A)}
$$
is normal on an interval $(a,b)$. Let $\{y_1,y_2,\dots,y_n\}$ be a
fundamental set of solutions of its complementary equation on $(a,b)$,
let $W$ be the Wronskian of $\{y_1,y_2,\dots,y_n\}$, and let $W_j$ be
the determinant obtained by deleting the last row and the $j$-th
column of $W$. Suppose $x_0$ is in $(a,b)$, let
$$
u_j(x)=(-1)^{(n-j)}\int_{x_0}^x{F(t)W_j(t)\over P_0(t)W(t)}\,dt,
\quad 1\le j\le n,
$$
and define
$$
y_p=u_1y_1+u_2y_2+\cdots+u_ny_n.
$$
\begin{alist}
\item % (a)
Show that $y_p$ is a solution of (A) and that
$$
y_p^{(r)}=u_1y^{(r)}_1+u_2y_2^{(r)}\cdots+u_ny^{(r)}_n,\quad 1
\le r \le n-1,
$$
and
$$
y_p^{(n)}=u_1y_1^{(n)}+u_2y_2^{(n)}+\cdots+u_ny_n^{(n)}+{F\over P_0}.
$$
\hint{See the derivation of the method of variation of
parameters at the beginning of the section.}
\item % (b)
Show that    $y_p$  is the
solution of the initial value problem
$$
\begin{array}{r}
P_0(x)y^{(n)}+P_1(x)y^{(n-1)}+\cdots+P_n(x)y=F(x),\\[5pt]
\hspace*{20pt} y(x_0)=0,\;
y'(x_0)=0,\dots,\quad y^{(n-1)}(x_0)=0.
\end{array}
$$
\item % (c)
Show that  $y_p$  can be written as
$$
y_p(x)=\int_{x_0}^x G(x,t)F(t)\,dt,
$$
where
$$
G(x,t)={1\over P_0(t)W(t)}\left|\begin{array}{cccc}
y_1(t)&y_2(t)&\cdots&y_n(t)\\[5pt]
y_1'(t)&y_2'(t)&\cdots&y_n'(t)\\[5pt]
\vdots&\vdots&\ddots&\vdots\\[5pt]
y_1^{(n-2)}(t)&y_2^{(n-2)}(t)&\cdots&y_n^{(n-2)}(t)\\[5pt]
y_1(x)&y_2(x)&\cdots&y_n(x)\end{array}\right|,
$$
which is called  the
\href{http://www-history.mcs.st-and.ac.uk/Mathematicians/Green.html}
{\color{blue}\it Green's function\/} for (A).
\item % (c)
Show that
$$
{\partial^{j}G(x,t)\over\partial x^j}
={1\over P_0(t)W(t)}\left|\begin{array}{cccc}
y_1(t)&y_2(t)&\cdots&y_n(t)\\[5pt]
y_1'(t)&y_2'(t)&\cdots&y_n'(t)\\[5pt]
\vdots&\vdots&\ddots&\vdots\\[5pt]
y_1^{(n-2)}(t)&y_2^{(n-2)}(t)&\cdots&y_n^{(n-2)}(t)\\[5pt]
y_1^{(j)}(x)&y_2^{(j)}(x)&\cdots&y_n^{(j)}(x)\end{array}\right|,\quad 0\le j\le
n.
$$
\item % (d)
Show that if $a<t<b$ then
$$
\left.{\partial^{j}G(x,t)\over\partial x^j}\right|_{x=t}=
\left\{\begin{array}{cl} 0,&1\le j\le n-2,\\[5pt]
\dst{1\over P_0(t)},&j=n-1.
\end{array}\right.
$$
\item % (e)
Show that
$$
 y_p^{(j)}(x)=\left\{\begin{array}{cl}\dst{\int_{x_0}^x
{\partial^{j}G(x,t)\over\partial x^j}F(t)\,dt},&1\le j\le n-1,\\[10pt]
\dst{{F(x)\over P_0(x)}+
\int_{x_0}^x{\partial^{(n)}G(x,t)\over\partial x^n}F(t)\,dt},&j=n.
\end{array}\right.
$$
\end{alist}

\exercisetext{In Exercises~\ref{exer:9.4.35}--\ref{exer:9.4.42}
use the method suggested by Exercise~\ref{exer:9.4.34} to find a particular
 solution in the form
$y_p=\int_{x_0}^x G(x,t)F(t)\,dt$, given the indicated fundamental
set of
 solutions.
 Assume that $x$ and $x_0$ are in an interval on which
the equation is normal.}

\item\label{exer:9.4.35}
$y'''+2y'-y'-2y=F(x); \quad \{e^x,\,e^{-x},e^{-2x}\}$

\item\label{exer:9.4.36}
$x^3y'''+x^2y''-2xy'+2y=F(x); \quad \{x,\,x^2,\,1/x\}$

\item\label{exer:9.4.37}
$x^3y'''-x^2(x+3)y''+2x(x+3)y'-2(x+3)y=F(x); \{x,x^2,xe^x\}$

\item\label{exer:9.4.38}
$x(1-x)y'''+(x^2-3x+3)y''+xy'-y=F(x); \quad \{x,\,1/x,\,e^x/x\}$

\item\label{exer:9.4.39}
$y^{(4)}-5y''+4y=F(x); \quad \{e^x,\,e^{-x},\,e^{2x},\,e^{-2x}\}$

\item\label{exer:9.4.40}
$xy^{(4)}+4y'''=F(x); \quad  \{1,\,x,\,x^2,\,1/x\}$

\item\label{exer:9.4.41}
$x^4y^{(4)}+6x^3y'''+2x^2y''-4xy'+4y=F(x)$;\quad
$\{x,x^2,1/x,1/x^2\}$

\item\label{exer:9.4.42}
$xy^{(4)}-y'''-4xy'+4y'=F(x); \quad \{1,\,x^2,\,e^{2x}, e^{-2x}\}$

\end{exerciselist}

